\color{unidarkgreen}{\subsection{Реле тока пуска охладителей}}
\color{black}

\begin{enumerate}[label=\arabic{section}.\arabic{subsection}.\arabic*,  align=left, labelsep=4pt, leftmargin=0pt, itemindent=57pt]

\item
Охлаждающие устройства в системах охлаждения мощных (авто)трансформаторов разбиты на
несколько групп, которые последовательно включаются в работу в зависимости от изменения условий
работы (авто)трансформатора. Функция пуска охладителей контролирует уровни тока на обмотках ВН,
НН, а также на общей обмотке автотрансформатора. 

\item
К каждой из обмоток подключено по 2 реле тока для пуска групп охладителей. Каждое реле имеет
индивидуально настраиваемый порог по току срабатывания, задаваемый пользователем.

\item
На рисунке \ref{pic:RTPO} представлена схема реле тока пуска охлаждения для обмотки ВН. Для общей обмотки автотрансформатора и обмотки НН схема аналогична.

\begin{figure}[H]
\centering
\includegraphics[width=0.8\textwidth,height=0.8\textheight,keepaspectratio]{\fbpath/041.1151 РТПО/img/RTPO.pdf}
\captionof{figure}{Функциональный блок РТПО}\label{pic:RTPO}
\end{figure}


\small
\begin{longtable}{|>{\centering\arraybackslash}m{5.3cm}|>{\centering\arraybackslash}m{4.2cm}|>{\centering\arraybackslash}m{1cm}|>{\centering\arraybackslash}m{3.3cm}|>{\centering\arraybackslash}m{1.8cm}|}
% Первая страница: полный заголовок
\caption{Перечень уставочных параметров функции пуска охладителей\hfill\vspace{-0.5\baselineskip}}\label{desc:tbl1}\\
\hline
\rowcolor{gray!30}
Наименование уставки & Диапазон & Шаг & Значение по умолчанию & Единица измерения \\
\hline
\endfirsthead

% Промежуточные страницы: «Продолжение таблицы»
\caption*{\hspace{3pt}\emph{Продолжение таблицы \ref{desc:tbl1}}\hfill\ }\vspace{-0.5\baselineskip} \\
\hline
\rowcolor{gray!30}
Наименование уставки & Диапазон & Шаг & Значение по умолчанию & Единица измерения \\
\hline
\endhead

%===m>ДЗАТ|-
%===m
\end{longtable}
\normalsize	

\end{enumerate}